% Main TFG LaTeX using esi-tfg class
\documentclass[spanish]{esi-tfg}

% Metadata and packages
% Metadata for TFG document - editable
\title{TFG: Gaia-X Salud — Dashboard de Pacientes}
\author{Enrique Mora Rodr\'iguez}
\director{F\'elix Jes\'us Villanueva Molina}
\degree{Grado en Ingeniería Informática}
\date{2026-04-30}

% Keywords
\keyWords{Gaia-X, FHIR, EDC, dashboard, microservicios}

% University / faculty metadata (edit as needed)
\faculty{Escuela Superior de Inform\'atica}
\university{Universidad de Castilla-La Mancha}
\city{Ciudad Real}

% If you prefer English, change documentclass option in main.tex


\begin{document}

\maketitle
\frontmatter
\begin{resumen}
% Resumen extraído del README — editar según convenga
Este Trabajo Fin de Grado implementa un MVP de un ecosistema de datos Gaia‑X en el dominio de la salud. Incluye tres microservicios (Provider, Trust, Consumer) y un dashboard React para visualización de datos clínicos FHIR, validación de políticas y trazabilidad end‑to‑end.
\end{resumen}

\tableofcontents
\mainmatter

\chapter{Introducción}
\section{Motivación}
% (editar) Breve motivación del proyecto

\section{Objetivos}
% Objetivos principales extraídos del README
\begin{itemize}
  \item Demostrar intercambio soberano de datos entre participantes
  \item Implementar un dashboard interactivo para datos FHIR
  \item Probar control de acceso y trazabilidad end‑to‑end
  \item Evaluar rendimiento bajo carga
\end{itemize}

\chapter{Estado del Arte}
\section{Gaia‑X y EDC}
\section{Estándares FHIR}
\section{Trabajos relacionados}

\chapter{Análisis y Diseño}
\section{Arquitectura del sistema}
% Insertar diagrama (por ejemplo: \includegraphics{figures/architecture.png})
\section{Modelado de datos}
\section{Decisiones arquitectónicas (ADR)}

\chapter{Implementación}
\section{Backend}
\subsection{Provider}
\subsection{Trust}
\subsection{Consumer}
\section{Frontend}
\section{Contenerización y despliegue}

\chapter{Pruebas y Validación}
\section{Pruebas unitarias e integrales}
\section{Ejecución E2E}
\section{Resultados de rendimiento}

\chapter{Conclusiones y Trabajo Futuro}
\section{Conclusiones}
\section{Líneas futuras}

\appendix
% Documentacion converted from documentacion.md
\chapter{Documentación y Base Bibliográfica}
\label{ch:documentacion}

Este capítulo recoge las fuentes primarias, estándares, especificaciones técnicas y referencias académicas en las que se fundamenta el diseño arquitectónico, tecnológico y funcional del proyecto. Las citas siguen el formato APA (editar según necesidad).

\section{Marco Conceptual: Gaia-X y Espacios de Datos}
Un \textbf{espacio de datos} (data space) es una infraestructura federada que permite el intercambio soberano de datos entre participantes sin centralización, garantizando que el propietario de los datos mantiene el control sobre quién los usa, cómo y para qué. Gaia-X es la iniciativa europea que define los principios y el marco.

\subsection{Documentos de referencia presentes en el repositorio}
\begin{itemize}
  \item \textbf{Whitepaper Gaia-X}: documentacion/WhitepaperGaiaX.pdf
  \item \textbf{Presentación de Casos de Uso en Salud}: documentacion/20240429-Presentaci\'on-Gaia-X-casos-de-uso-SALUD.pdf
\end{itemize}

\section{Eclipse Dataspace Connector (EDC)}
Se explica el control plane y data plane y la estrategia EDC-first adoptada en el proyecto (ADR-002). Referencias: Eclipse EDC docs y Dataspace Protocol.

\section{Estándar HL7 FHIR R4}
Breve descripción de FHIR, recursos usados (Bundle, Observation) y ejemplo de Observation para presión arterial.

\section{Sistemas de codificación: LOINC y UCUM}
Listado de códigos LOINC usados (55284-4, 85354-9, 8480-6, 8462-4, 8867-4) y unidades UCUM (`mm[Hg]`, `/min`).

\section{Regulación: GDPR y EHDS}
Resumen de principios GDPR relevantes (minimización, finalidad, integridad) y alineamiento con EHDS.

\section{Seguridad: OAuth2, JWT y Control de Acceso}
Descripción del uso del flujo client\_credentials y JWT en el Trust Service; evolución prevista a DID/VC.

\section{Arquitectura de Microservicios y DDD}
Principios aplicados (SRP, API First, deploy independiente) y estructura de carpetas (api/, domain/).

\section{Stacks Tecnológicos}
Resumen de stacks backend y frontend (Spring Boot + Java 25, React + TypeScript, Vite, TanStack Query, Recharts, Tailwind).

\section{Contenedores, DevOps y Observabilidad}
Multi-stage builds, health checks, Docker Compose y uso de Actuator/Micrometer/OpenTelemetry.

\section{Testing y Calidad}
Pir\'amide de tests, frameworks usados (JUnit 5, Vitest, RTL) y criterios de cobertura.

\section{Documentación interna del proyecto}
Listado de documentos internos y su ruta: plan/analisis_funcional.md, plan/analisis_tecnico.md, plan/contratos_api.md, plan/tareas.md, plan/evidencias_t9.md, plan/informe_carga_t10.md, plan/METADATA.md.

\section{ADRs y Bibliografía}
Resumen de ADRs (ADR-001..ADR-007) y referencia a la bibliografía completa. Para la bibliografía vea el archivo documentacion.md o references.bib.

% Fin de documentacion.tex

% Defensa técnica convertido desde TFG_explicacion.md
\chapter{Explicación Técnica para la Defensa}
\label{ch:defensa}

\section{Visión general del proyecto}
Resumen del problema, objetivos del MVP y casos de uso clínicos implementados (presión arterial y frecuencia cardíaca). El sistema opera con más de 10.000 eventos clínicos de prueba.

\section{Arquitectura general}
Diagrama de sistema (ver figura \ref{fig:architecture}) y flujo E2E: Publicación → Solicitud → Autorización → Consumo → Visualización.

\section{Provider Node (Puerto 8081)}
Rol: publicador de datos FHIR. APIs destacadas: POST /api/v1/datasets/fhir y GET /api/v1/datasets. Se validan Bundles y Observations; modelo de datos DatasetRecord.

\section{Trust Service (Puerto 8082)}
Rol: árbitro y Authorization Server. Evalúa políticas, emite JWT (client_credentials) y registra audit trail.

\section{Consumer Node (Puerto 8083)}
Rol: solicita autorizaciones, ejecuta jobs de consumo, expone API para el dashboard y almacena resultados en memoria (MVP) con opción a PostgreSQL en producción.

\section{Frontend: Dashboard de Pacientes}
Componentes: SearchBar, MetricsLineChart, DistributionPie, Table. Configuración Vite + TanStack Query + Recharts.

\section{Comunicación entre servicios}
HTTP REST entre servicios; headers: Authorization: Bearer <JWT>, X-Request-Id para correlaci\'on; GAIAX_TRUST_BASE_URL configurada por entorno.

\section{Testing y Resultados de Rendimiento}
Resumen de estrategia de tests y KPIs obtenidos (Error rate 0.00%, P95 latency 23-27ms, Throughput 150-168 req/s). Referencias a plan/informe_carga_t10.md.

\section{Despliegue}
Instrucciones rápidas: docker compose -f compose.yaml up -d --build; healthchecks y dependencias por readiness.

\section{Preguntas frecuentes del tribunal}
Sección con preguntas y respuestas preparadas (p.ej., por qué FHIR, por qué EDC-first, cómo se asegura la soberanía de datos, limitaciones del MVP).

% Fin de defensa.tex


\chapter{Instrucciones de desarrollo}
% Copiar fragmentos relevantes del README o referenciarlo

\chapter{Evidencias}
% Añadir logs, capturas y resultados de carga

\bibliographystyle{plain}
\bibliography{references}

\end{document}
